\section{Almost sure invariance}

The extreme decomposition rests on a statement about a single Markov kernel on a measurable space
$(\Omega,\mathcal F)$; it is proved at that generality.

\begin{definition}[The a.s.\ invariant $\sigma$-algebra, Georgii (7.3)]
    \label{def:ext-ae-invariant}
    \lean{ProbabilityTheory.Kernel.aeInvariantSets,
      ProbabilityTheory.Kernel.aeInvariantSigmaAlgebra,
      ProbabilityTheory.Kernel.aeInvariantSigmaAlgebra_le}
    \leanok

    Let $\pi$ be a Markov kernel from $(\Omega,\mathcal F)$ to itself and $\mu$ a finite measure
    with $\mu\pi = \mu$. The system
    \[\mathcal I_\pi(\mu) \;=\; \set{A \in \mathcal F \;:\; \pi(A \mid \cdot) = \1_A \ \mu\text{-a.s.}}\]
    of \textbf{$\mu$-almost surely $\pi$-invariant} sets is a sub-$\sigma$-algebra of $\mathcal F$.
\end{definition}
\begin{proof}
    \leanok

    Stability under complements is the Markov property. For a countable union $A = \bigcup_n A_n$
    of members, $\pi(A \mid \cdot) \geq \pi(A_n \mid \cdot) = \1_{A_n}$ a.s.\ for each $n$, so
    $\1_A \leq \pi(A \mid \cdot)$ a.s.; since $\int \pi(A\mid\cdot)\,\dd\mu = \mu\pi(A) = \mu(A) =
    \int \1_A \,\dd\mu < \infty$, the inequality is an equality a.s.
\end{proof}

\begin{theorem}[Densities of invariant measures, Georgii (7.3)]
    \label{thm:ext-invariant-density}
    \uses{def:ext-ae-invariant}
    \lean{ProbabilityTheory.Kernel.invariant_withDensity_of_measurable,
      ProbabilityTheory.Kernel.measurable_of_invariant_withDensity,
      ProbabilityTheory.Kernel.invariant_withDensity_iff_measurable}
    \leanok

    Let $\pi$, $\mu$ be as in Definition \ref{def:ext-ae-invariant}. For measurable
    $f : \Omega \to [0,\infty]$ with $\int f \,\dd\mu < \infty$,
    \[(f\mu)\pi = f\mu \quad\Longleftrightarrow\quad f \text{ is } \mathcal I_\pi(\mu)\text{-measurable}.\]
\end{theorem}
\begin{proof}
    \leanok

    ($\Leftarrow$) By monotone approximation it suffices to treat $f = c\1_A$ with
    $A \in \mathcal I_\pi(\mu)$, and there $\int_A \pi(B \mid \cdot)\,\dd\mu = \mu(A \cap B)$
    because $\pi$ neither leaves $A$ nor enters it.

    ($\Rightarrow$) Fix $0 < c < \infty$ and put $A = \set{f \geq c}$. Invariance of $f\mu$ and of
    $\mu$ give two conservation identities; combined with $(f - c)(1 - \pi(A\mid\cdot)) \geq 0$
    on $A$ they yield
    $c\int_{A^c} \pi(A\mid\cdot) \leq \int_{A^c} f\,\pi(A\mid\cdot) \leq c\int_{A^c}\pi(A\mid\cdot)$.
    Since $f < c$ on $A^c$, equality forces $\pi(A \mid \cdot) = 0$ a.s.\ on $A^c$, and the
    conservation identity for $\mu$ then forces $\pi(A\mid\cdot) = 1$ a.s.\ on $A$. So
    $A \in \mathcal I_\pi(\mu)$.
\end{proof}

\begin{theorem}[Extremality, Georgii (7.4)]
    \label{thm:ext-74}
    \uses{thm:ext-invariant-density}
    \lean{ProbabilityTheory.Kernel.aeInvariantSigmaAlgebraFamily,
      ProbabilityTheory.Kernel.eq_of_absolutelyContinuous_of_trivialOn,
      ProbabilityTheory.Kernel.mem_extremePoints_iff_trivialOn}
    \leanok

    Let $\Pi$ be a non-empty family of Markov kernels on $(\Omega,\mathcal F)$ and $\mathcal P_\Pi$
    the convex set of probability measures invariant under every member. For $\mu \in \mathcal
    P_\Pi$ put $\mathcal I_\Pi(\mu) = \bigcap_{\pi \in \Pi} \mathcal I_\pi(\mu)$. Then $\mu$ is
    extreme in $\mathcal P_\Pi$ if and only if $\mu$ is trivial on $\mathcal I_\Pi(\mu)$.
\end{theorem}
\begin{proof}
    \leanok

    If $A \in \mathcal I_\Pi(\mu)$ has $0 < \mu(A) < 1$ then $\mu(\cdot \mid A)$ and
    $\mu(\cdot \mid A^c)$ lie in $\mathcal P_\Pi$ by Theorem \ref{thm:ext-invariant-density}, and
    they split $\mu$. Conversely, if $\mu = a\nu_1 + b\nu_2$ with $a,b>0$ then $\nu_1 \ll \mu$, so
    $\nu_1 = f\mu$ for $f = \dd\nu_1/\dd\mu$; Theorem \ref{thm:ext-invariant-density} makes $f$
    measurable for $\mathcal I_\Pi(\mu)$, triviality makes it a.e.\ constant, and its integral $1$
    makes that constant $1$.
\end{proof}

\begin{example}[Almost sure, not strict, Georgii (7.5)]
    \label{ex:ext-75}
    \uses{thm:ext-74}
    \lean{ProbabilityTheory.Kernel.Example75.ker,
      ProbabilityTheory.Kernel.Example75.mu,
      ProbabilityTheory.Kernel.Example75.invariant_mu,
      ProbabilityTheory.Kernel.Example75.eq_empty_or_univ_of_strictly_invariant,
      ProbabilityTheory.Kernel.Example75.trivial_on_strictly_invariant,
      ProbabilityTheory.Kernel.Example75.not_mem_extremePoints}
    \leanok

    In Theorem \ref{thm:ext-74}, $\mathcal I_\Pi(\mu)$ cannot be replaced by the $\sigma$-algebra of
    \emph{strictly} invariant sets $\set{A : \pi(A \mid \cdot) = \1_A \text{ everywhere}}$. Take
    $\Omega = \set{0,1,2}$ and
    \[\pi \;=\; \begin{pmatrix} 1 & 0 & 0 \\ \tfrac12 & 0 & \tfrac12 \\ 0 & 0 & 1\end{pmatrix},
      \qquad \mu \;=\; \tfrac12\delta_0 + \tfrac12\delta_2 .\]
    Then $\mu$ is invariant, and a strictly invariant $A$ satisfies
    $\tfrac12\1_A(0) + \tfrac12\1_A(2) = \1_A(1)$, which forces $A = \emptyset$ or $A = \Omega$;
    so $\mu$ is trivial on the strictly invariant sets. But $\delta_0$ and $\delta_2$ are
    invariant, so $\mu$ is not extreme.
\end{example}

\begin{proposition}[The deterministic case is Mathlib's ergodicity]
    \label{prop:ext-mathlib}
    \uses{thm:ext-74, ex:ext-75}
    \lean{ProbabilityTheory.Kernel.invariant_deterministic_iff,
      ProbabilityTheory.Kernel.mem_aeInvariantSets_deterministic_iff,
      ProbabilityTheory.Kernel.preErgodic_iff_trivialOn_aeInvariant}
    \leanok

    Let $T : \Omega \to \Omega$ be measurable and $\pi = \delta_{T(\cdot)}$. Then $\mu\pi = \mu$
    if and only if $T$ preserves $\mu$, and $\mathcal I_\pi(\mu)$ consists of the measurable $A$
    with $T^{-1}A =_\mu A$. For measure-preserving $T$ on a probability space, the zero-one law over
    the \emph{strictly} invariant sets $T^{-1}A = A$ is equivalent to triviality on
    $\mathcal I_\pi(\mu)$; Theorem \ref{thm:ext-74} thus specialises to the characterisation of the
    ergodic measures as the extreme points of the invariant ones.

    The equivalence corrects an a.e.\ invariant set to a strictly invariant one by iterating $T$,
    which uses measurability of $T$; by Example \ref{ex:ext-75} no such correction exists for a
    general kernel.
\end{proposition}

\begin{corollary}[Pre-ergodicity is an almost-everywhere notion]
    \label{cor:ext-preergodic-ae}
    \uses{prop:ext-mathlib}
    \lean{ProbabilityTheory.Kernel.aeInvariantSets_congr,
      ProbabilityTheory.Kernel.preErgodic_iff_forall_nullMeasurableSet,
      ProbabilityTheory.Kernel.preErgodic_congr_ae}
    \leanok

    For measure-preserving $T$ on a probability space, the zero-one law over the strictly invariant
    sets is equivalent to the zero-one law over the null measurable $A$ with $T^{-1}A =_\mu A$; in
    particular it depends only on the $\mu$-a.e.\ class of $T$. Likewise $\mathcal I_\pi(\mu)$
    depends only on the $\mu$-a.e.\ class of $\pi$.
\end{corollary}

\section{Tail triviality and asymptotic independence}

\begin{theorem}[Georgii (7.9)]
    \label{thm:ext-79}
    \lean{MeasureTheory.forall_measure_eq_zero_or_one_iInf_iff,
      MeasureTheory.forall_measure_eq_zero_or_one_iInf_iff_of_directed,
      MeasureTheory.GibbsMeasure.forall_tail_measure_eq_zero_or_one_iff,
      MeasureTheory.GibbsMeasure.isTailTrivial_iff_asymptotically_independent}
    \leanok

    Let $S$ be countable and $\mu \in \Prob(\Cfg, \mathcal F)$. The following are equivalent.
    \begin{enumerate}
        \item $\mu$ is trivial on $\Tail$.
        \item For every event $A$,
          \[\lim_{\Lambda \in \Fin} \ \sup_{B \in \mathcal F_{S \setminus \Lambda}}
            \abs{\mu(A \cap B) - \mu(A)\mu(B)} \;=\; 0 .\]
    \end{enumerate}
    The same equivalence holds on any probability space for an antitone family
    $(\mathcal A_i)_{i \in I}$ of sub-$\sigma$-algebras indexed by a countable directed non-empty
    $I$, with $\mathcal F_{S\setminus\Lambda}$ replaced by $\mathcal A_i$ and $\Tail$ by
    $\bigcap_{i} \mathcal A_i$.
\end{theorem}
\begin{proof}
    \leanok

    ($1 \Rightarrow 2$) By L\'evy's downward theorem $\mu[\1_A \mid \mathcal F_{S \setminus \Lambda}]
    \to \mu[\1_A \mid \Tail]$ in $L^1$, and triviality on $\Tail$ makes the limit the constant
    $\mu(A)$. For $B \in \mathcal F_{S\setminus\Lambda}$,
    \[\mu(A \cap B) - \mu(A)\mu(B)
      = \int_B \bigl(\mu[\1_A \mid \mathcal F_{S\setminus\Lambda}] - \mu(A)\bigr) \dd\mu,\]
    which is bounded in absolute value by the $L^1$ distance, hence tends to $0$.

    ($2 \Rightarrow 1$) Any $A \in \Tail$ lies in every $\mathcal F_{S\setminus\Lambda}$, so taking
    $B = A$ gives $\abs{\mu(A) - \mu(A)^2} \leq \varepsilon$ for every $\varepsilon > 0$, whence
    $\mu(A) \in \set{0,1}$.

    The directed case reduces to the sequential one along a monotone cofinal sequence in $I$.
\end{proof}

\section{Uniform mixing}

\begin{definition}[Uniform mixing, Georgii (7.10)]
    \label{def:ext-710}
    \lean{MeasureTheory.GibbsMeasure.IsUniformlyMixing}
    \leanok

    A probability measure $\mu$ on $\Cfg$ is \textbf{uniformly mixing} if for every local event
    $A$ and every $\varepsilon > 0$, all sufficiently large $\Lambda \in \Fin$ satisfy
    $|\mu(A \mid B) - \mu(A)| \le \varepsilon$ for every non-null
    $B \in \Tail_\Lambda = \mathcal{F}_{\Lambda^{\mathrm{c}}}$.
\end{definition}

\begin{proposition}[Georgii (7.11)(a)]
    \label{prop:ext-711a}
    \uses{def:ext-710, def:gibbs-meas}
    \lean{MeasureTheory.GibbsMeasure.eq_of_isGibbsMeasure_lambdaSpecification_of_isUniformlyMixing,
      MeasureTheory.GibbsMeasure.G_lambdaSpecification_eq_singleton_of_isUniformlyMixing}
    \leanok

    Let $\gamma = \rho\lambda$ be a quasilocal $\lambda$-specification with everywhere positive
    densities over a $\sigma$-finite non-zero a priori measure, and let $\mu \in \Gibbs{\gamma}$ be
    uniformly mixing. Then $\Gibbs{\gamma} = \set{\mu}$.
\end{proposition}
\begin{proof}
    \uses{lem:remark-128-2}
    \leanok

    Let $\nu \in \Gibbs{\gamma}$, $A$ a local event and $\varepsilon > 0$. Pick $\Lambda \in \Fin$
    from uniform mixing for $A$ and $\varepsilon$. Since $\gamma_\Lambda(A \mid \cdot)$ is a
    version of $\mu[\1_A \mid \Tail_\Lambda]$, the bound on $\mu(A \mid B)$ over non-null
    $B \in \Tail_\Lambda$ gives $|\gamma_\Lambda(A \mid \cdot) - \mu(A)| \le \varepsilon$
    $\mu$-a.s. Quasilocality provides $\Delta \in \Fin$ and an $\mathcal F_\Delta$-measurable $f$
    with $\sup_\eta |\gamma_\Lambda(A \mid \eta) - f(\eta)| \le \varepsilon$, so
    $\set{|f - \mu(A)| > 2\varepsilon} \in \mathcal F_\Delta$ is $\mu$-null, hence $\nu$-null by
    Lemma \ref{lem:remark-128-2}, and $|\gamma_\Lambda(A \mid \cdot) - \mu(A)| \le 3\varepsilon$
    $\nu$-a.s. Integrating $\nu\gamma_\Lambda = \nu$ gives $|\nu(A) - \mu(A)| \le 3\varepsilon$.
    Hence $\nu$ and $\mu$ agree on the generating $\pi$-system of local events, so $\nu = \mu$.
\end{proof}

\begin{proposition}[Georgii (7.11)(b)]
    \label{prop:ext-711b}
    \uses{def:ext-710, prop:ext-711a, def:gibbs-meas}
    \lean{MeasureTheory.GibbsMeasure.eventually_forall_abs_real_apply_sub_le_of_forall_isGibbsMeasure_eq,
      MeasureTheory.GibbsMeasure.tendsto_iSup_ofReal_abs_real_apply_sub_of_forall_isGibbsMeasure_eq,
      MeasureTheory.GibbsMeasure.tendsto_finiteVolumeDistributions_of_forall_mem_GP_eq,
      MeasureTheory.GibbsMeasure.tendsto_iSup_ofReal_abs_integral_sub_of_forall_isGibbsMeasure_eq,
      MeasureTheory.GibbsMeasure.isUniformlyMixing_of_forall_isGibbsMeasure_eq,
      MeasureTheory.GibbsMeasure.eventually_forall_abs_real_apply_sub_le_lambdaSpecification_of_G_eq_singleton,
      MeasureTheory.GibbsMeasure.tendsto_iSup_ofReal_abs_integral_sub_lambdaSpecification_of_G_eq_singleton,
      MeasureTheory.GibbsMeasure.G_lambdaSpecification_eq_singleton_iff_isUniformlyMixing,
      MeasureTheory.GibbsMeasure.G_lambdaSpecification_eq_singleton_iff_isUniformlyMixing_of_isFiniteMeasure}
    \leanok

    Let $E$ be standard Borel and $\gamma$ a quasilocal specification such that for every
    $\Lambda \in \Fin$ there is a finite measure $\kappa_\Lambda$ with
    $\gamma_\Lambda(A \mid \omega) \le \kappa_\Lambda(A)$ for all $\omega \in \Cfg$ and all
    $A \in \mathcal F_\Lambda$. Let $\mu \in \Prob(\Cfg)$ be such that every Gibbs measure of
    $\gamma$ equals $\mu$ (in particular, $\Gibbs{\gamma} = \set{\mu}$). Then for every local
    event $A$ and every $\varepsilon > 0$, all sufficiently large $\Lambda \in \Fin$ satisfy
    $|\gamma_\Lambda(A \mid \omega) - \mu(A)| \le \varepsilon$ simultaneously for all
    $\omega \in \Cfg$; hence $\sup_\omega |\gamma_\Lambda(A \mid \omega) - \mu(A)| \to 0$,
    $\sup_\omega |\gamma_\Lambda(f \mid \omega) - \mu(f)| \to 0$ for every bounded local
    observable $f \in \mathcal{L}$, and $\gamma_\Lambda(\cdot \mid \omega) \to \mu$ in the
    topology of local convergence for every $\omega$. If moreover $\mu \in \Gibbs{\gamma}$, then
    $\mu$ is uniformly mixing.

    The domination hypothesis holds for a $\lambda$-specification $\gamma = \rho\lambda$ with
    volumewise bounded densities over a probability a priori measure $\lambda$. For such $\gamma$
    with everywhere positive densities, and for a finite non-zero $\lambda$ with densities bounded
    relative to $\lambda/\lambda(E)$, $\Gibbs{\gamma} = \set{\mu}$ if and only if
    $\mu \in \Gibbs{\gamma}$ is uniformly mixing.
\end{proposition}
\begin{proof}
    \uses{thm:exg-4-17}
    \leanok

    Were the first conclusion to fail, some local $A$, some $\varepsilon > 0$ and boundary
    conditions $\omega_\Lambda$ along a subnet would keep
    $|\gamma_\Lambda(A \mid \omega_\Lambda) - \mu(A)| > \varepsilon$. The domination hypothesis
    makes the net $(\gamma_\Lambda(\cdot \mid \omega_\Lambda))_\Lambda$ locally equicontinuous, so
    it has a cluster point, which is a Gibbs measure by Theorem \ref{thm:exg-4-17} because
    $\gamma$ is quasilocal, hence equals $\mu$ --- contradicting the separation on $A$. A bounded
    local observable is uniformly approximated by finite linear combinations of indicators of
    local events, which gives the statement for $f \in \mathcal L$. Mixing follows: conditioning
    on a non-null $B \in \Tail_\Lambda$ averages the numbers $\gamma_\Lambda(A \mid \omega)$, all
    within $\varepsilon$ of $\mu(A)$. The equivalence combines this with Proposition
    \ref{prop:ext-711a}.
\end{proof}
